\documentclass[10pt,a4paper]{amsart}
\usepackage[margin=19mm]{geometry}
\usepackage{amsmath,amssymb,mathtools}
\usepackage[scr=rsfs,cal=euler]{mathalfa}
\usepackage{xfrac}
\usepackage[unicode,hidelinks,bookmarks=false]{hyperref}
\newcommand{\ie}{i.e.\ }
\newcommand{\cf}{cf.\ }
\newcommand{\resp}{respectively\ }
\newcommand{\Subsection}{\subsection}
\providecommand{\nobreakdash}{\nobreak\hbox{-}}
\setlength{\emergencystretch}{2em}
\multlinegap0pt
\usepackage{bm}
\usepackage{bbm}
\usepackage{dsfont}
\usepackage{xspace}
\usepackage{cancel}
\usepackage{mathtools}
\usepackage{amsthm}
\makeatletter
\@ifundefined{theorem}{\newtheorem{theorem}{Theorem}}{}
\@ifundefined{Lemma}{\newtheorem{Lemma}[theorem]{Lemma}}{}
\@ifundefined{proposition}{\newtheorem{proposition}[theorem]{Proposition}}{}
\makeatother
\newcommand{\ind}{\mathrm{ind}}
\newcommand{\op}[1]{\operatorname{#1}}
\title[On Malle's conjecture for the product of symmetric and nilpo]{On Malle's conjecture for the product of symmetric and nilpotent groups}
\author{Hrishabh Mishra, Anwesh Ray}
\date{Source: arXiv:2402.01189v3 (CC BY 4.0)}
\begin{document}
\maketitle

\noindent\emph{Source and licence.} On Malle's conjecture for the product of symmetric and nilpotent groups. Version arXiv:2402.01189v3 (CC BY 4.0), distributed under the Creative Commons Attribution 4.0 International licence, as recorded in the arXiv licence metadata for this exact version and shown on the abstract page https://arxiv.org/abs/2402.01189v3. Full rights notice: licenses/arxiv-2402.01189-CC-BY-4.0.txt.\par\smallskip

\noindent\textit{Editorial scope.} Counted unit: the Proposition whose source label is \texttt{ind(sigma, g) propn} in the article (source0.tex lines 347--364, source label \texttt{ind(sigma, g) propn}), delivered together with the article's own complete original proof of it (source0.tex lines 365--381, one \texttt{proof} environment of 17 lines). No mathematics is written, repaired or re-derived: statement and proof are the article's own text line for line. Required context retained uncounted: index\_comp (lines 333--338). Every definition, display and label of those lines is retained unchanged. Omitted from this package and not redistributed: 1--332, 339--346, 382--1071. Nothing inside a retained excerpt is omitted or altered: every retained body is delivered exactly as the article's lines are written, so no omission receipt of any kind accompanies this package. Citations: the retained excerpts carry no citation command at all, so no bibliography entry of the article is transcribed and no reference is redistributed. Reference adaptation: none was needed and none was made; every reference inside the delivered unit resolves inside it, so no unresolved reference or citation is produced. Printed slips kept literal: no printed word, symbol, hypothesis or number of the counted statement or of its proof was altered. Layout and class: the article's own class is \texttt{amsart} with packages including geometry, amscd, amsmath, amsxtra, amsthm, amssymb, stmaryrd, xr, mathrsfs, mathtools, enumerate, commath, comment, xcolor; the wrapper is the lane's pinned \texttt{amsart} editorial wrapper at 10pt on A4, which restates the theorem-like environments the retained text needs and adds the article's own macro definitions that the retained text needs (\texttt{\textbackslash ind}, \texttt{\textbackslash op}), with extra wrapper packages (bm, bbm, dsfont, xspace, cancel, mathtools). The delivered numbering is the unit's own. Assets: no retained span contains a figure environment, a graphics-inclusion command, a PDF-page inclusion, a table or a TikZ picture; no image file is copied into this package, the package declares no asset dependency and no image file is redistributed with it. Licence: version arXiv:2402.01189v3 of this article is distributed under the Creative Commons Attribution 4.0 International licence, as recorded in the arXiv licence metadata for this exact version and shown on the abstract page https://arxiv.org/abs/2402.01189v3. A full rights notice is delivered at licenses/arxiv-2402.01189-CC-BY-4.0.txt. Characters: every retained excerpt is pure ASCII, so the delivered engine, XeLaTeX with its default Latin Modern text font, reports no missing character.
\begin{Lemma}\label{index_comp}
    Let $\sigma \in S_n$ and $\tau \in S_m$ with disjoint cycle decomposition $\sigma = \prod_i c_j$ and $\tau = \prod_j d_j$. Suppose that $(|c_i|,|d_j|)= 1$ for all $i,j$ and consider $(\sigma, \tau) \in S_{n}\times S_m \subset S_{nm}$ then
    \[
\ind(\sigma, \tau) = \ind(\sigma)\cdot m+ \ind(\tau)\cdot n- \ind(\sigma) \cdot \ind(\tau).
\]
\end{Lemma}

\begin{proposition}\label{ind(sigma, g) propn}
    Let $G$ be a finite nilpotent group in its regular representation $G\hookrightarrow S_{|G|}$.
    \begin{enumerate}
        \item\label{c1 ind(sigma, g) propn} When $n=3$, suppose that $2 \nmid |G|$. Then we have that \begin{equation}\label{ind(sigma, g) propn f1}
        \ind((12),g)/|G| > 2 \text{ and } \ind((123),g)/|G| > 1
        \end{equation}
        for all $g \in G$.
        \item\label{c2 ind(sigma, g) propn} For $n=4$ assume that $2,3 \nmid |G|$. Then \begin{equation}\label{ind(sigma, g) propn f2}
        \ind((12),g)/|G| > 2\text{ and } \ind((12)(34),g)/|G| > 1, 
        \end{equation}
        \begin{equation}\label{ind(sigma, g) propn f3}
        \ind((123),g)/|G| > 3\text{ and } \ind((1234),g)/|G| > 2
        \end{equation}
        for all $g \in G$.
        \item\label{c3 ind(sigma, g) propn} When $n=5$ assume that $2,3,5 \nmid |G|$. Then we have that \begin{equation}\label{ind(sigma, g) propn f4}
        \ind(\sigma,g)/|G| \geq 1 + \ind(\sigma)-1/7\end{equation} for all $\sigma \in S_5$ and $g \in G$. Further if $\sigma$ is not conjugate to $(12345)$ in $S_5$ then \begin{equation}\label{ind(sigma, g) propn f5}\ind(\sigma,g)/|G| \geq 12/7 + \ind(\sigma).\end{equation}
    \end{enumerate}
\end{proposition}

\begin{proof} Since $g$ sits inside $S_m$ via the regular representation, we have that \begin{equation}\label{ind g equation}\op{ind}(g)=|G|\left(1-\frac{1}{e(g)}\right)\end{equation} where $e(g)$ denotes the order of $g$. Note that since it is assumed that $\ell_G>2$, and hence $e(g)>2$ for all $g\neq 1_G$. It follows that $\op{ind}(g)>\frac{|G|}{2}$ for all $g\in G^\ast$ and \[\op{ind}(G)=|G|\left(1-\frac{1}{\ell_G}\right).\]
\par We begin with \eqref{c1 ind(sigma, g) propn}. For $(12) \in S_3$, using Lemma \ref{index_comp} we deduce that for all $g \in G$ we have that $\ind((12),g) = |G| + 2\cdot \ind(g)$ as $\ind(12) = 1$. We deduce that \[\ind((12),g)/|G| \geq 1 + 2\cdot\ind(G)/|G| > 2\]
since $\ell_G > 2$.
\par For $(123) \in S_3$, if $3 \nmid |G|$, then again using Lemma \ref{index_comp} we obtain that for all $g \in G$, $\ind((123),g) = 2|G| + \ind(g)$ as $\ind(123) = 2$. Hence the inequality holds in this case. If $3 \mid |G|$, then $G= G(3)\times \widetilde{G}$. Given any $g \in G$ we can write it as $(g(3), \widetilde{g}) \in S_{|G(3)|}\times S_{\widetilde{G}} \subset S_{|G|}$. Further we note that $\ind((123),g) = \ind((123),(g(3),\widetilde{g})) = \ind(((123),g(3)),\widetilde{g})$ in $S_{|G|}$. Since $3 \nmid |\widetilde{G}|$ we can use Lemma \ref{index_comp} to deduce that \[\ind(((123),g(3)),\widetilde{g}) = \ind((123,g(3))\cdot(|\widetilde{G}|-\ind(\widetilde{g}))+3\cdot|G(3)|\cdot\ind(\widetilde{g}).\] Now, since the smallest possible value of $\ind((123,g(3))$ is $2|G(3)|$ when $g(3)$ is identity and smallest possible value for $\ind(\widetilde{g})$ is $\ind(\widetilde{G})$ by definition. Hence, \[
\ind((123,g(3))\cdot(|\widetilde{G}|-\ind(\widetilde{g}))+3\cdot|G(3)|\cdot\ind(\widetilde{g}) \geq 2|G| + |G(3)|\cdot\ind(\widetilde{G}) > |G|.\] This concludes the proof.
\par \textit{(2)} We note that  $\ind(12),$ $\ind(123),$ $\ind(1234),$ and $\ind(12)(34)$ are $1, 2, 3,$ and $2$ respectively. Further $2,3 \nmid |G|$ so we can use Lemma \ref{index_comp} and we deduce that for all $g \in G$, $\ind((12),g) = |G| + 3\cdot\ind(g) > 2|G|$, $\ind((123),g) = 2|G| + 2\cdot\ind(g) > 3|G|$ as $\ind(g) \geq \ind(G) = \frac{\ell_G-1}{\ell_G} |G| \geq \frac{4}{5}|G|$, next $\ind(1234, g) = 3|G| + \ind(g) > 2|G|$ and $\ind((12)(34), g) = 2|G| + 2\cdot \ind(g) > |G|$.
\par \textit{(3)} Using Lemma \ref{index_comp} we conclude that 
\[\begin{split}
    \ind(\sigma, g) = & \ind(\sigma)|G| + 5\cdot\ind(g) - \ind(\sigma)\cdot\ind(g) \\
   \geq  & \ind(\sigma)\cdot|G| + (5-\ind(\sigma))\cdot\ind(g) \geq \ind(\sigma)\cdot |G| + (5-\ind(\sigma))\cdot \ind(G) \\
   \geq  & \ind(\sigma)\cdot |G| + (5-\ind(\sigma))\cdot (\ell_G-1)|G|/\ell_G \\
   \geq & \ind(\sigma)\cdot |G| + (5-\ind(\sigma))\cdot 6|G|/7 \\
   \geq & \op{ind}(\sigma) |G|+6|G|/7,
\end{split} \] since $\ell_G \geq 7$. Dividing both sides of the above inequality by $|G|$, we find that 
\[\frac{\op{ind}(\sigma, g)}{|G|}\geq 1+\op{ind}(\sigma) -\frac{1}{7},\] and this proves \eqref{ind(sigma, g) propn f4}. 
For $\sigma \in S_5$ not conjugate to $(12345)$ we have that $(5-\ind(\sigma))\geq 2$ and hence \[\ind(\sigma)\cdot |G| + (5-\ind(\sigma))\cdot 6|G|/7 \geq \ind(\sigma)\cdot |G| + 12|G|/7,\] from which \eqref{ind(sigma, g) propn f5} follows. This concludes the proof of Proposition.
\end{proof}

\end{document}
